\documentclass[11pt]{ctexbook}
\usepackage[margin=1in]{geometry}
\usepackage{hyperref}
\usepackage{enumitem}
\hypersetup{colorlinks=true, linkcolor=blue, urlcolor=blue}
\setlist[itemize]{leftmargin=1.6em}

\title{Berkeley Agentic AI Fall 2025\\对 Spring 2025 主书的课程级增补}
\author{Supplement Workspace: berkeley\_agentic\_ai\_f25}
\date{2026-04-16}

\begin{document}
\maketitle
\tableofcontents

\chapter{Berkeley Agentic AI Fall 2025：对 Spring 2025 主书的课程级增补}
\label{chap:berkeley_agentic_ai_f25_extension}

\section{本章范围与证据边界}

这一章不是重新把 \emph{CS294/194-196: Agentic AI (Fall 2025)} 逐讲完整转写一遍，而是基于官方公开来源，给 \emph{CS294/194-280: Advanced Large Language Model Agents (Spring 2025)} 主书补上一层课程级增量解释。证据边界只有四类：

\begin{itemize}
\item 官方课程页：\url{https://rdi.berkeley.edu/agentic-ai/f25}
\item 官方 MOOC 页：\url{https://agenticai-learning.org/f25}
\item 官方 Berkeley RDI playlist：\\ \url{https://www.youtube.com/playlist?list=PLS01nW3RtgoqGkm4UeqNeZLccW-OGc1fJ}
\item 官方课程页显式列出的 slides、readings，以及 10 月 6 日那条官方但未列入 playlist 的 unlisted recording：\\ \url{https://www.youtube.com/watch?v=VfOA2a0dj4w}
\end{itemize}

因此，本章的立场非常明确：

\begin{itemize}
\item 能从 official page、slides、playlist、official readings 直接支撑的内容，就按事实写。
\item 只能从 lecture 标题和 reading 标题推断的内容，要明确说这是保守解释，而不是已验证的 slide-level 细节。
\item 本章服务于 Spring 2025 主书的升级，而不是替代主书。
\end{itemize}

这门 Fall 2025 课的一个关键信号来自 MOOC 页本身。官方直接写出：它建立在 Fall 2024 LLM Agents MOOC 和 Spring 2025 Advanced LLM Agents MOOC 的基础之上。这意味着 Fall 2025 不是简单重复，而是在原来 reasoning、planning、tools、multimodal、theorem proving、safety 的框架上，再往评测、可验证后训练、multi-agent、scientific discovery、deployment、embodiment 这些更系统工程的方向推进。

\section{为什么 Fall 2025 是主书的真实后续}

Spring 2025 主书已经覆盖了高级 LLM agent 的核心研究主线：推理时计算、学习推理、记忆与规划、开放 post-training recipe、coding agents 与漏洞检测、多模态环境、formal mathematics，以及安全与权限边界。

如果只看这些章节，会以为 Berkeley 这条线已经把 agent 讲完了。但 Fall 2025 的官方课程描述和具体 syllabus 显示，它把重点向另外几个方向明显移动：

\begin{itemize}
\item evaluation 从附属 benchmark 变成课程主轴；
\item post-training 不再只是 preference optimization，而是 environment-feedback-aligned、verifier-centered 的 agent training；
\item agent 的真正难点被重新定位到 system design、infra、grader、orchestration、observability、access control、simulation；
\item multi-agent 和 recursive self-improvement 被提升为独立主题；
\item scientific discovery、paper agents、embodiment 把 agent 从会用工具的语言模型推向可在复杂环境中持续运行的研究或控制系统。
\end{itemize}

换句话说，Spring 2025 更像是在回答：高级 agent 需要哪些能力模块；而 Fall 2025 更像是在追问：当这些能力真的要落地到 benchmark、训练 recipe、产品部署、群体交互和现实环境时，系统应该如何组织。

\section{从 LLM 能做什么，转向 agent system 如何组织}

9 月 15 日的 \emph{LLM Agents Overview} slides 给了一个非常关键的框架。公开 slide 文本里明确出现：

\begin{itemize}
\item Pretraining
\item Reasoning RL
\item Classic post-training / RLHF
\item Evaluation
\item Systems and infra to scale
\end{itemize}

这件事的重要性不在于它重新讲一遍 LLM pipeline，而在于它把 evaluation 和 systems \& infra 提升到了与模型算法并列的位置。Spring 2025 主书当然已经涉及环境反馈、tool use 和系统闭环，但 Fall 2025 在课程一开始就告诉你：agent 时代的关键瓶颈不再只是模型会不会推理，而是训练、评测、部署、扩展这整条链路能不能稳定运行。

9 月 22 日的 \emph{Evolution of system designs from an AI engineer perspective} 又把这个方向推进了一步。Yangqing Jia 的公开 slides 里最值得注意的，不是某个具体模型，而是三个判断：

\begin{itemize}
\item 新算法会持续推动模型能力；
\item 应用空间会持续膨胀；
\item \texttt{AI infra has become the 3rd pillar of enterprise IT strategy}。
\end{itemize}

这说明 Fall 2025 对 agent 的看法已经明显偏向系统工程对象。主书在 Spring 2025 里更多讨论模型、verifier、tool loop、workflow design；而 Fall 2025 在这里开始强调 developer efficiency、infra efficiency、multi-cloud supply chain management，以及把 development、training 和 inference 统一到 AI-native platform。

\section{从偏好对齐后训练走向可验证 agent 训练}

\subsection{Jiantao Jiao：agentic model 不再只是 human preference model}

9 月 29 日的 \emph{Post-Training Verifiable Agents} slides 里，最关键的定义不是某个 loss，而是对模型目标的重新改写：

\begin{itemize}
\item Earlier Chat Models: Human Aligned Models
\item Agentic Models: Environment Feedback Aligned Models
\end{itemize}

这个区分非常重要。它意味着系统的优化目标，从让用户偏好更满意，转向让 agent 在环境中取得可验证的正确结果。Jiantao 的 slides 明确把 agent training 拆成三个要素：

\begin{itemize}
\item environment
\item tools
\item verifier
\end{itemize}

也就是说，一个 agentic model 不是只在上下文窗口里生成看起来合理的文本；它要读取环境状态、调用工具、再根据 verifier 返回的信号被训练。这里的 verifier 可以是 code unit tests、proof checker、DOM script、ground-truth state verifier。

\subsection{Weizhu Chen：真正困难的是 grader、rubric 和 product constraints}

如果 Jiantao 讲的是目标函数变了，那么 10 月 13 日的 \emph{Some Challenges and Lessons from Training Agentic Models} 讲的是：真正难的是怎么把这个目标做成能训练、能上线、能稳定优化的工程对象。

公开 slides 里有一个很典型的 coding-agent 结构：problem prompt、code repository、tool calls / tool responses、Linux system、unit tests、model patch + PR description。这说明 agent training 已经不是对着答案做监督学习，而是要在一个类似软件工程 sandbox 的环境里学习一整套交互行为。Weizhu 的 slides 还反复强调三件事：

\begin{itemize}
\item rubric 必须可评分、可解释、可组合；
\item verifiable data 与 non-verifiable data 需要混合治理；
\item product grader 远比单一 unit test 复杂。
\end{itemize}

对于软件 agent，公开 slides 明确列出 grader 的多个层次：unit-test grader、patch grader、rollout grader、behavior grader、user-experience grader、task-specific grader。主书里我们已经能看到 coding agent、漏洞检测、formal verification、OS 环境等内容，但 Fall 2025 让一个更工业级的事实变得清楚：agent 训练的难点，经常不是模型结构，而是你能否定义出一套足够细、足够可靠、又不会把系统带偏的 grader stack。

\subsection{主书应如何升级对后训练的理解}

如果把 Spring 2025 和 Fall 2025 放在一起看，一个更完整的 post-training 图景是：

\begin{itemize}
\item preference alignment 解决“用户觉得像不像一个好助手”
\item reasoning RL 解决“模型会不会在有明确答案的问题上形成更强推理轨迹”
\item verifiable agent training 解决“模型能不能在环境里持续地产生可被 checker、tests、rubric 验证的动作和结果”
\end{itemize}

因此，主书里有关 post-training 的章节，在读完这个 supplement 后，最好升级成一句更严格的话：post-training 的最终对象不是静态答案，而是受环境、工具、verifier、grader 和 product constraints 共同约束的 agent policy。

\section{evaluation 不再是结果汇报，而是 harness 本身}

\subsection{Oct 6：先建立 taxonomy，再谈 benchmark}

10 月 6 日的 \emph{Agent Evaluation \& Project Overview} 虽然只是课程中的一讲，但它在整体结构里非常关键。公开 slides 明确给出四组区分：

\begin{itemize}
\item close-ended vs.\ open-ended
\item verifiable vs.\ non-verifiable
\item static vs.\ dynamic
\item taxonomy of agent eval tasks
\end{itemize}

同时 slides 明确提出了 Outcome Validity。一个 agent eval 不是只要能自动评分就行，而是要问：这个任务真的对应我们想测的 capability 吗？这个评分信号能区分真正完成任务与投机性通过吗？这个环境能否逼迫 agent 暴露 memory、planning、tool use、adaptation 等真实能力？

\subsection{Survey：agent eval 已经形成独立知识结构}

官方 reading \emph{Survey on Evaluation of LLM-based Agents} 把 agent eval 拆成四层：

\begin{itemize}
\item capability evaluation：planning、tool use、self-reflection、memory
\item application-specific evaluation：web、software engineering、scientific、conversational agents
\item generalist agent evaluation
\item evaluation frameworks
\end{itemize}

这说明 Fall 2025 相比 Spring 2025 的一个重要变化是：课程不再把 evaluation 看成某个 domain 的局部细节，而是把它视为组织整个 agent field 的骨架。

\subsection{Oct 27：没有误差条的 leaderboard 很可能会误导你}

10 月 27 日的 \emph{Predictable Noise in LLM} 与官方 reading \emph{Adding Error Bars to Evals} 一起，把另一个经常被忽略的问题推到了台前：小 benchmark、hard benchmark、generative benchmark 的统计噪声很大。

公开 slides 与 reading 共同支持几个关键判断：

\begin{itemize}
\item 很多 agent eval 的 test set 比 ImageNet 小得多；
\item 但每道题的生成和执行成本更高；
\item 只看 pass rate 的百分点差异，经常无法判断提升是否稳健；
\item 对模型比较应该更多使用 paired comparison，而不是只比较 summary statistic；
\item evaluation 设计需要 power analysis，而不是做完再看数字。
\end{itemize}

所以，在新的 agent textbook 里，evaluation 不能只写成列举 benchmark。它必须被写成 task design、verifier design、statistical design、reporting discipline 四者共同组成的 harness。

\section{multi-agent 不只是分工协作，而是 population-level dynamics}

\subsection{Noam Brown：self-play、minimax 与 exploitability}

10 月 20 日的公开 slides 直接从 self-play、minimax equilibrium、population best response、exploitability 出发。这里真正值得带回主书的不是扑克例子本身，而是下面这条思想：当系统变成多主体交互时，能力的定义不再只是单个 agent 的平均任务成功率，而是它在一个 population 里的稳定性、可被利用程度，以及能否通过 self-play 或 recursive improvement 继续优化。

这会改变我们对 agent scaling 的理解。很多人把 scaling 想成更多参数、更多 inference-time compute、更多工具。但 Fall 2025 的 multi-agent 视角说明，另一个可扩展维度是：让 agent 在相互博弈、协作或竞争的环境中，学习更强的策略稳定性。

\subsection{Nov 17：LLM 时代的 multi-agent system 更接近系统课题}

11 月 17 日官方只公开了题目 \emph{Multi-Agent Systems in the Era of LLMs} 和录播，没有公开 slides 或 readings。因此这里必须保守。最稳妥的解释是：课程在 10 月 20 日的 self-play / game-theoretic view 之后，又安排了一讲更系统导向的 multi-agent session，说明 multi-agent 已经不只是理论点缀，而是 Berkeley 认为值得单独展开的主轴。

\section{agent 开始进入 scientific discovery，并把论文本身变成 agent}

Spring 2025 主书第 11 讲已经谈到 abstraction and discovery with large language model agents。但 Fall 2025 把 discovery 进一步推向了更具体、更工程化的对象。

11 月 3 日的题目是 \emph{AI Agents to Automate Scientific Discoveries}。官方 reading 一篇是 Nature 文章 \emph{The Virtual Lab of AI agents designs new SARS-CoV-2 nanobodies}，另一篇是 \emph{Paper2Agent: Reimagining Research Papers As Interactive and Reliable AI Agents}。

\emph{Paper2Agent} 的公开摘要尤其关键。它明确提出：系统可以用多个 agents 去分析论文和代码库，构造 Model Context Protocol server，再通过自动生成和运行测试来提高可靠性。这不是简单的论文问答机器人，而是在把 research artifact 变成带工具接口、能执行 workflow、能被验证的 agent interface。

因此，读完这组内容之后，主书里的 research assistant 不应再被想成一个会 summarize paper 的模型，而应被想成一个更强的对象：它理解论文、它能调用论文配套代码、它能重放或校验原文结果、它能在新 query 下执行可追溯 workflow。

\section{deployment 才是 agent 真正进入现实世界时最厚的一层}

11 月 10 日的 \emph{Practical Lessons from Deploying Real-World AI Agents} 是 Fall 2025 最像现实世界 agent 课的一讲。Clay Bavor 的公开 slides 里最值得带回主书的是 \emph{Agent Iceberg}。

表层当然还是熟悉的东西：LLM、RAG、tool use。但真正厚重的底层是：

\begin{itemize}
\item complex workflows \& orchestration
\item observability and monitoring
\item regression testing
\item user simulation
\item prompt injection protection
\item knowledge partitioning
\item role-based access controls
\item reporting \& audit
\item staging and release management
\item model migration and failover
\end{itemize}

这张图的意义很大。它把 Spring 2025 主书中许多分散出现的安全、tooling、workflow、benchmark、memory、access control 线索，压缩成一个非常工程化的结论：真正可部署的 agent，主要成本不在“会不会答题”，而在底层平台能不能持续地测试、观察、隔离、上线、回滚、审计和模拟。

\section{从多模态和 GUI，走向 embodied / autonomous agents}

Spring 2025 主书已经有两讲很强的 multimodal material：\emph{Multimodal Autonomous AI Agents} 和 \emph{Multimodal Agents -- From Perception to Action}。这些章节已经把 perception-action loop、GUI grounding、OS/Web benchmarks 讲得很细。Fall 2025 的 12 月 1 日 \emph{Autonomous Agents: Embodiment, Interaction, and Learning} 再往前走了一步：它不只是讨论屏幕上的交互，而是把 agent 放回 embodied / autonomous settings。

虽然官方没有公开 slides，但课程提供的两个 readings 已经把方向讲清楚了：

\begin{itemize}
\item \emph{Outracing Champion Gran Turismo Drivers with Deep Reinforcement Learning}
\item \emph{SLAC: Simulation-Pretrained Latent Action Space for Whole-Body Real-World RL}
\end{itemize}

这说明课程在这里想要强调的，不是简单“机器人也能用 LLM”，而是：

\begin{itemize}
\item embodied agent 的动作空间更连续、更受物理约束；
\item 训练往往更依赖 simulation、latent action abstraction、control pipeline；
\item interaction loop 从浏览器和 GUI 的离散操作，转向与真实世界状态耦合的连续决策。
\end{itemize}

\section{安全与安全性不再只是尾声提醒}

12 月 8 日的 closing session 叫 \emph{Agentic AI Safety \& Security}。公开材料只有课程页和官方录播，没有公开 slides 或 readings，因此这里只能保守解释。但它在课程结构里的位置本身就很有信息量：它是整门课的收官讲次，而且直接把 Safety 和 Security 放进标题。

这说明 Fall 2025 对安全问题的态度，比“顺便提醒一下风险”更强。它更像在说：当 agent 真的进入 training、evaluation、deployment、scientific workflow、multi-agent interaction 和 embodied setting 之后，安全与安全性必须作为 system-level closure 被重新处理。

\section{主书最应该如何重构}

如果要把 Spring 2025 主书升级成更贴近 2025 年底 agent 版图的 textbook，至少应做五个结构性调整：

\begin{itemize}
\item 把 evaluation 提前成全书前置框架，而不是章末 benchmark 列表；
\item 把 post-training 重写成 verifier、grader、environment、rubric 的系统问题；
\item 给 deployment 单独留出一个厚章节；
\item 把 multi-agent 从拓展话题升成独立 part；
\item 给 scientific discovery 与 embodied agents 各自留出更明确的结构位置。
\end{itemize}

\section{本章小结}

如果只用一句话概括 Berkeley \emph{Agentic AI (Fall 2025)} 对 Spring 2025 主书的增量，那就是：它把 advanced LLM agents 从能力模块的集合，进一步推进成需要训练、评测、部署、治理和群体交互的系统工程对象。

更具体地说，Fall 2025 至少完成了八个推进：

\begin{itemize}
\item 把 evaluation 从结果汇报提升为 harness 主轴；
\item 把 post-training 从 preference optimization 推向 verifier-centered agent training；
\item 把 system design / infra 推到 agent 讨论中心；
\item 把 multi-agent 提升为 population-level dynamics 问题；
\item 把 scientific discovery 推向可执行的 paper agents / virtual labs；
\item 把 deployment 写成可观测、可测试、可审计的产品栈；
\item 把 multimodal 继续推进到 embodied / autonomous agents；
\item 把 safety \& security 作为整门课的收束，而不是尾声装饰。
\end{itemize}

\section{局限与后续工作}

本 supplement 是扎实的课程级补章，但仍然有三个残余风险：

\begin{itemize}
\item Sep 8、Nov 3、Nov 17、Dec 1、Dec 8 没有完整的公开 slides 或 transcript 层，因此这些部分必须比 Spring 2025 主书更保守；
\item Sep 29 的两篇官方 OpenAI readings 在当前环境下只能稳定拿到标题和 URL，不能稳定抽取正文；
\item 本 workspace 只做 course-level extension，没有为 Fall 2025 每一讲重建新的 transcript.jsonl / coverage\_units.jsonl / lecture.tex 流水线。
\end{itemize}

因此，本章最适合的角色是：作为 Spring 2025 主书的高质量后续课程增量章。如果后续要把 Fall 2025 完整升格成第二本教材，就需要再为每一讲单独跑 lecture-level harness。

\end{document}
